¿Crees que tus métodos preventivos son suficientes?

\textbf{Respuesta:}
No, no son 100\% suficientes, pues, como vimos en clase, en la ciberseguridad no existe la seguridad absoluta y considerarse totalmente a salvo es uno de los mayores errores que podemos cometer.

Como puede verse, la seguridad no es un producto ni un estado estático, sino un proceso continuo. Si bien las medidas preventivas cotidianas que utilizo logran mitigar la gran mayoría de los ataques oportunistas y automatizados, existen vectores de riesgo que escapan de mi control.

\begin{itemizeColor}
    \item \textbf{Vulnerabilidades de día cero (\textit{Zero-Days}):} Fallos de seguridad en el sistema operativo, controladores o navegadores que son desconocidos para los desarrolladores y para los cuales aún no existe un parche disponible.
    \item \textbf{Ataques a la cadena de suministro (\textit{Supply Chain Attacks}):} Situaciones en las que un atacante compromete una biblioteca de software, actualización legítima o dependencia antes de que llegue a los usuarios finales.
    \item \textbf{Filtraciones en servidores de terceros:} Aunque un usuario mantenga contraseñas únicas y robustas, la información personal sensible almacenada en servidores de empresas, universidades o entidades gubernamentales puede verse expuesta si dichas plataformas sufren una brecha de seguridad.
    \item \textbf{Acciones ajenas a mí en sistemas compartidos:} Dado que vivo con mi familia y me conecto a redes públicas, como la red de la facultad, no están en mis manos las acciones descuidadas que puedan tomar los miembros de mi familia o mis compañeros por el desconocimiento del peligro de sus acciones, exponiendo mi seguridad aunque tome las acciones de protección que mencioné antes.
\end{itemizeColor}

Por ello, los estándares actuales de seguridad, tales como la Arquitectura de Confianza Cero (\textit{Zero Trust}) del NIST~\cite{nist_zero_trust}, promueven el principio de \textbf{asumir la brecha} (\textit{Assume Breach}). Esto significa que las medidas preventivas son indispensables para elevar el costo del ataque, pero deben complementarse siempre con una actitud vigilante, capacidad de detección temprana, contención de daños y políticas periódicas de respaldo y recuperación.

